\documentclass[11pt,letterpaper]{article}
% Reusable style. Build each .tex from its own directory using scripts/build.sh.
\usepackage[T1]{fontenc}
\usepackage[utf8]{inputenc}
\usepackage{lmodern}
\usepackage[margin=0.86in,headheight=15pt]{geometry}
\usepackage{amsmath,amssymb,amsthm,mathtools}
\usepackage{booktabs,longtable,tabularx,array,multicol}
\usepackage[table]{xcolor}
\usepackage{microtype}
\usepackage{enumitem}
\usepackage{fancyhdr}
\usepackage{titlesec}
\usepackage{hyperref}
\usepackage{listings}
\usepackage[most]{tcolorbox}
\usepackage{needspace}
\definecolor{aeNavy}{HTML}{172D43}
\definecolor{aeBlue}{HTML}{176B90}
\definecolor{aeTeal}{HTML}{148477}
\definecolor{aePale}{HTML}{F2F7F9}
\definecolor{aeGray}{HTML}{5C6771}
\hypersetup{colorlinks=true,linkcolor=aeBlue,urlcolor=aeBlue,pdfauthor={Algorithm Engineering Documentation Suite},pdfsubject={Independent educational materials based on the supplied TAOCP volume topics}}
\titleformat{\section}{\Large\bfseries\color{aeNavy}}{\thesection}{0.6em}{}
\titleformat{\subsection}{\large\bfseries\color{aeBlue}}{\thesubsection}{0.6em}{}
\titleformat{\subsubsection}{\normalsize\bfseries\color{aeNavy}}{\thesubsubsection}{0.6em}{}
\setlist[itemize]{leftmargin=*,topsep=3pt,itemsep=1pt}
\setlist[enumerate]{leftmargin=*,topsep=3pt,itemsep=3pt}
\setlength{\parindent}{0pt}\setlength{\parskip}{5pt}
\pagestyle{fancy}\fancyhf{}\fancyhead[L]{\footnotesize\color{aeGray}ALGORITHM ENGINEERING}\fancyhead[R]{\footnotesize\color{aeGray}DOCUMENTATION SUITE}\fancyfoot[L]{\footnotesize\color{aeGray}Independent study and engineering practice}\fancyfoot[R]{\footnotesize\thepage}
\renewcommand{\headrulewidth}{0.35pt}
\newtcolorbox{engineeringnote}[1][]{breakable,colback=aePale,colframe=aeBlue,boxrule=0.5pt,arc=1mm,fonttitle=\bfseries,title={Engineering note},#1}
\newtcolorbox{evidencebox}[1][]{breakable,colback=white,colframe=aeTeal,boxrule=0.6pt,arc=1mm,fonttitle=\bfseries,title={Evidence and deliverables},#1}
\lstdefinestyle{algo}{basicstyle=\ttfamily\footnotesize,breaklines=true,frame=single,rulecolor=\color{aeGray!40},backgroundcolor=\color{aePale},showstringspaces=false,tabsize=2,columns=fullflexible,keepspaces=true,keywordstyle=\color{aeBlue}\bfseries,commentstyle=\color{aeTeal}}
\lstset{style=algo}
\newcommand{\DocTitle}[3]{%
\hypersetup{pageanchor=false}\begin{titlepage}\vspace*{1.2cm}
{\color{aeTeal}\rule{\linewidth}{2.3pt}}\par\vspace{1.15cm}
{\fontsize{24}{30}\selectfont\bfseries\color{aeNavy}\raggedright #1\par}\vspace{0.45cm}
{\Large\color{aeBlue}#2\par}\vspace{1.3cm}
{\large\bfseries #3\par}\vspace{0.35cm}
{\color{aeGray}Edition 1.0 \quad | \quad 29 October 2026\par}\vfill
{\color{aeGray}\small Original educational material aligned to the chapter and section topics in the five supplied volumes of Donald E. Knuth's \emph{The Art of Computer Programming}. This document is not an authorized reproduction or substitute for the original books.}\par\vspace{0.4cm}
{\color{aeTeal}\rule{\linewidth}{0.75pt}}
\end{titlepage}\hypersetup{pageanchor=true}\pagenumbering{roman}}
\newcommand{\SourceNote}[1]{\par\smallskip{\footnotesize\color{aeGray}\textbf{Source locator:} #1}\par\smallskip}
\newcommand{\Outcome}[1]{\par\textbf{Competency target.} #1\par}
\newcommand{\Checkpoint}[1]{\par\textbf{Checkpoint.} #1\par}
% Intentionally simple printable writing lines: call \AnswerLines{number}.
\newcount\alcount
\newcommand{\AnswerLines}[1]{\par\alcount=0\loop\ifnum\alcount<#1\par\noindent\rule{\linewidth}{0.3pt}\vspace{0.33cm}\advance\alcount by1\repeat\par}
\newcommand{\FormField}[2][0.62\linewidth]{\noindent\textbf{#2}\hspace{0.65em}\rule{#1}{0.4pt}\par\vspace{0.18cm}}
\newcommand{\ReviewTable}{\begin{tabularx}{\linewidth}{@{}p{0.23\linewidth}X p{0.2\linewidth}@{}}\toprule Criterion & Recorded evidence & Rating (0--4)\\\midrule Correctness & & \\[1em] Proof quality & & \\[1em] Complexity & & \\[1em] Testing & & \\[1em] Reproducibility & & \\[1em] Communication & & \\[1em]\bottomrule\end{tabularx}}

\begin{document}
\DocTitle{Algorithm Verification and Review Procedure}{Correctness, numerical reliability, security and performance gates}{Engineering documentation}
\tableofcontents
\clearpage
\pagenumbering{arabic}
\section{Review intent}
Treat algorithm selection as a justified design decision. The procedure applies to educational prototypes, reusable libraries, and performance-critical services. Strength of evidence should scale with risk: an input parser and cryptographic token generator require security scrutiny beyond ordinary benchmark correctness.
\section{Gate 1: specification and assumptions}
Write domains, preconditions, output semantics, mutation, determinism and numerical representation. Define behavior for invalid inputs, empty inputs, duplicate keys and resource exhaustion. List assumptions whose violation would invalidate the proof.
\section{Gate 2: invariant and progress review}
Inspect initialization, preservation, exit condition and termination. For recursive procedures, define base cases and decrease toward them. For backtracking, prove that state restoration is exact after each branch, including exceptions and early returns.
\section{Gate 3: test engineering}
Apply equivalence classes, boundary cases, negative tests, randomized property checks, cross-language differential tests and exhaustive small-instance enumeration. For a SAT solver, a satisfying witness is checked clause by clause; an UNSAT claim from a complex solver may require a separately verifiable proof certificate or small exhaustive oracle.
\section{Gate 4: numerical and resource safety}
Inspect overflow/underflow, type conversions, precision, untrusted input sizes, allocations, stack depth, recursion limits, algorithmic complexity attacks, nontermination and key-dependent branching. Explicitly prohibit simulation-grade PRNGs for security-critical token generation.
\section{Gate 5: complexity and workload selection}
Identify theoretical worst case, typical workload, domain restrictions, memory use, and data structure operations. Do not approve claims of guaranteed constant time if hash collisions or rehashing produce worse cases. Compare alternatives over a workload matrix.
\section{Gate 6: experiment and decision}
Follow ALG-BENCH-01; archive measurements and counterexamples. A decision note must state chosen algorithm, rejected alternatives, costs accepted, validation evidence and conditions under which the choice should be revisited.
\section{Review outcome rubric}
\begin{tabularx}{\linewidth}{@{}p{0.22\linewidth}X@{}}\toprule Status & Decision rule\\\midrule
Approved & All required gates satisfied; remaining limitations clearly documented and accepted.\\
Conditionally approved & Specific non-critical evidence gaps have owner, scope and follow-up deadline.\\
Rejected & Incorrectness, missing termination reasoning, inappropriate randomness, unsafe arithmetic or irreproducible claims remain unresolved.\\
\bottomrule\end{tabularx}
\SourceNote{Vol.~1, Chapters 1--2; Vol.~2, Chapters 3--4; Vol.~3, Chapters 5--6; Vol.~4B, \S7.2.2. Source volumes provide algorithms and analysis; review workflow is independently authored.}
\end{document}
